
\begin{table}[H]
\centering
\caption{Summary Statistics}
\begin{threeparttable}
\begin{tabular}{lcccc}
\toprule
Variable & Mean & SD & Min & Max \\
\midrule
\multicolumn{5}{l}{\textit{Panel A: Treatment}} \\
Immigrant share, 2008 (\%) & 5.86 & 3.06 & 2.67 & 17.96 \\[0.5em]
\multicolumn{5}{l}{\textit{Panel B: Second-generation outcomes}} \\
Employment rate, 2022 (\%) & 72.1 & 11.8 & 0.0 & 87.5 \\
Tertiary education share, 2023 (\%) & 17.4 & 7.0 & 0.0 & 39.0 \\
Primary-only education share (\%) & 22.7 & 7.2 & 0.0 & 54.0 \\
Employment gap vs.\ Danish origin (pp) & 10.0 & 10.8 & -7.9 & 77.2 \\[0.5em]
\multicolumn{5}{l}{\textit{Panel C: Municipality characteristics}} \\
Total population (2008) & 109,422 & 249,091 & 6,712 & 1,645,825 \\
Non-Western descendant count & 2,386 & 7,283 & 20 & 62,142 \\
Danish-origin employment rate (\%) & 82.2 & 3.2 & & \\
\bottomrule
\end{tabular}
\begin{tablenotes}[flushleft]
\small
\item \textit{Notes:} N = 100 municipalities. The sample includes all Danish municipalities (post-2007 boundaries) with at least 50 non-Western descendants in 2008. Immigrant share is the number of non-Western immigrants divided by total population in 2008Q1 (FOLK1C). Employment rates are from Statistics Denmark RAS200 for non-Western descendants aged 25--39. Tertiary education includes short-cycle higher education, bachelor, and master/PhD programs (HFUDD11). Employment gap is the difference between Danish-origin and non-Western descendant employment rates.
\end{tablenotes}
\end{threeparttable}
\label{tab:summary}
\end{table}
